% Boyer実験ガイド
\documentclass[12pt,a4paper]{article}

\usepackage{fontspec}
\usepackage{xeCJK}
\setCJKmainfont{Hiragino Mincho ProN}
\setCJKsansfont{Hiragino Kaku Gothic ProN}
\usepackage[margin=2.5cm]{geometry}
\usepackage{amsmath,amssymb}
\usepackage{hyperref}
\usepackage{tcolorbox}
\usepackage{booktabs}
\usepackage{fancyhdr}

\tcbuselibrary{breakable,skins}

\newtcolorbox{storybox}[1][]{colback=blue!5,colframe=blue!60!black,
  fonttitle=\bfseries\sffamily,title={#1},breakable,arc=3mm}
\newtcolorbox{mathbox}[1][]{colback=green!5,colframe=green!50!black,
  fonttitle=\bfseries\sffamily,title={#1},breakable,arc=3mm}
\newtcolorbox{resultbox}[1][]{colback=yellow!8,colframe=orange!80!black,
  fonttitle=\bfseries\sffamily,title={#1},breakable,arc=3mm}
\newtcolorbox{honestbox}[1][]{colback=gray!8,colframe=gray!50!black,
  fonttitle=\bfseries\sffamily,title={#1},breakable,arc=3mm}
\newtcolorbox{expbox}[1][]{colback=red!5,colframe=red!60!black,
  fonttitle=\bfseries\sffamily,title={#1},breakable,arc=3mm}

\setlength{\parskip}{0.8em}
\setlength{\parindent}{0pt}

\pagestyle{fancy}
\fancyhf{}
\rhead{Wright Brothers}
\lhead{Boyer実験ガイド}
\cfoot{\thepage}

\begin{document}

\begin{center}
{\Huge\bfseries 50年越しの実験}\\[0.5cm]
{\Large Boyerの予測を初めて検証する}\\[0.3cm]
{\Large 素数が真空を変える}\\[1cm]
{\large 初心者のための理論と実験ガイド}\\[0.5cm]
{\large Wright Brothers Research Program, 2026年4月}
\end{center}

\vspace{0.5cm}
\tableofcontents
\newpage

%====================================================================
\part{理論: なぜ真空に力があるか}
%====================================================================

\section{カシミール効果（1948年）}

\begin{storybox}[何もない空間に力がある]
2枚の金属板をナノメートルの距離で向かい合わせると、板の間に引力が生じる。電荷もなく、磁石でもないのに。

原因: 量子力学では「完全に何もない」ことが不可能。真空は常に微小な電磁場の揺らぎ（ゼロ点揺動）で満たされている。板を置くと揺らぎのパターンが変わり、板の間のエネルギーが外より低くなる。エネルギーが低い方に引かれる → 引力。

1948年にカシミールが理論的に予測。1997年に精密測定で確認。
\end{storybox}

\section{揺らぎの「倍音」}

\begin{mathbox}[モードと倍音]
板の間の電磁場には、ギターの弦と同じように倍音（モード）がある:
$$f_1, \; 2f_1, \; 3f_1, \; 4f_1, \; 5f_1, \ldots$$

各モードがゼロ点エネルギー $\hbar\omega/2$ を持つ。全モードの合計が真空エネルギー。

合計: $1 + 2 + 3 + 4 + 5 + \cdots$（1Dの場合）

1D（パイプ）なら $\zeta(-1) = 1+2+3+\cdots = -1/12$。

しかし実際の実験は3D（平行板）。3Dでは横方向の自由度が加わり、和は $\sum n^3$ になる:
$$\zeta(-3) = 1^3 + 2^3 + 3^3 + 4^3 + \cdots = +\frac{1}{120}$$

カシミールエネルギー密度: $\rho = -\pi^2\hbar c/(720 a^4)$。
マイナス符号は3Dの積分から来る。$720 = 6/\zeta(-3)$。

結果: 板の間のエネルギーが外より\textbf{低い} → \textbf{引力}。

これがカシミール効果の数学的正体。実験で精度1\%以内で確認済み。
\end{mathbox}

\section{Boyerの予測（1974年）}

\begin{resultbox}[もし片端を変えたら？]
Boyer は考えた: 片方の板を「完全磁気導体（PMC）」に変えたらどうなるか？

通常の板（PEC）: 電場がゼロになる（ディリクレ境界条件）
PMC: 磁場がゼロになる（ノイマン境界条件）

PEC-PMC にすると、残るモードは奇数だけ:
$$1, \; 3, \; 5, \; 7, \; 9, \ldots$$

3Dでの合計（$\sum n^3$ の奇数版）:
$$\zeta_{\neg 2}(-3) = 1^3 + 3^3 + 5^3 + 7^3 + \cdots = -\frac{7}{120}$$

3Dカシミールエネルギー密度は $\rho \propto -\zeta(-3)$ なので:
$$\rho_{DN} \propto -\zeta_{\neg 2}(-3) = -(-7/120) = +\frac{7}{120} > 0$$

\textbf{正！ 符号が反転した！ しかも大きさは通常の7倍！}

板の間のエネルギーが外より高い → \textbf{斥力}。板が反発する。

（1Dでも同様: $\zeta(-1) = -1/12 \to \zeta_{\neg 2}(-1) = +1/12$ で反転。）
\end{resultbox}

\begin{honestbox}[50年間誰もこれを実験で確認していない]
理由: PMC（完全磁気導体）が存在しない。

普通の金属は「完全電気導体」に近い。しかし「完全磁気導体」として振る舞う物質は、自然界にもメタマテリアルでも実現されていない（カシミールに関わる光の周波数帯では）。

Boyer の予測は理論的に正しいと証明されている（Kenneth et al. 2002, PRL）。しかし実験で確認されたことは一度もない。
\end{honestbox}

\section{「カシミール斥力」を名乗る実験たち}

\begin{honestbox}[なぜ既存の実験はBoyerの斥力ではないか]
\textbf{Munday et al. (2009, Nature)}: 金とシリカの間にブロモベンゼン液を満たし、カシミール斥力を測定。しかしこれは「液体の誘電率が中間的だから」斥力が生じる（DLP条件）。板の間の\textbf{モード構造は変わっていない}。全モードが存在したまま。素数ミュートではない。

\textbf{Tang et al. (2017, Nature Photonics)}: T字型シリコンナノ構造をスライドさせ、力のベクトル和が正になる配置を発見。しかし\textbf{全ペアのカシミール力は引力のまま}。複数の引力の合計が偶然正の方向を向いただけ。モード構造は変わっていない。

\textbf{Boyer 球殻 (1968)}: 導体球殻のカシミール自己エネルギーが正（外向き圧力）。しかし球のモードは全スペクトル（球面調和関数の全次数）。偶数を消していない。

\textbf{メタマテリアル提案 (Rosa 2008, Leonhardt 2007)}: メタマテリアルでPMCを近似する理論提案。しかしKramers-Kronig制約により、狭帯域の磁気応答は虚数周波数で消える。実現していない。

\textbf{全てのケースで共通}: モード和は $\sum n^{-s}$（全モード）のまま。偶数モードが消えていない。オイラー積は無傷。

\textbf{モード和が $1^3+3^3+5^3+\cdots = \zeta_{\neg 2}(-3)$ に切断された実験は、人類史上ゼロ。}
\end{honestbox}

%====================================================================
\part{再解釈: 素数がカギだった}
%====================================================================

\section{オイラー積}

\begin{mathbox}[ゼータ関数の「素因数分解」]
リーマンゼータ関数は素数で分解できる（オイラー、1737年）:
$$\zeta(s) = \frac{1}{1-2^{-s}} \times \frac{1}{1-3^{-s}} \times \frac{1}{1-5^{-s}} \times \cdots$$

各因子が1つの素数に対応。
\end{mathbox}

\section{素数2を「消す」}

\begin{mathbox}[素数ミュート]
$p=2$ の因子を除去する:
$$\zeta_{\neg 2}(s) = \zeta(s) \times (1 - 2^{-s}) = \sum_{n \text{ odd}} n^{-s}$$

$2$ の因子を消すと、$2$ で割れる全ての整数（偶数）が和から消え、奇数だけが残る。

3Dでは:
$$\zeta(-3) = +\frac{1}{120} \;\xrightarrow{\text{素数2ミュート}}\; \zeta_{\neg 2}(-3) = -\frac{7}{120}$$

Boyer の計算はまさにこれ: 「偶数モードが消えて奇数モードだけ残る」= オイラー積の $p=2$ 因子の除去。

\textbf{Boyer は「素数2をミュートした」ことに気づいていなかった。}
\end{mathbox}

\section{PMC ではなく「幾何学」}

\begin{resultbox}[問題の変換]
\textbf{Boyer の問い}: PMC を見つけよ → 50年間未解決

\textbf{我々の問い}: 奇数モードだけを選ぶ幾何学を見つけよ → \textbf{$\lambda/4$ 共振器}

$\lambda/4$ 共振器 = 片端を閉じて片端を開けた管。ビール瓶と同じ。
\begin{itemize}
\item $\lambda/2$（両端閉じ）: 全倍音 $1, 2, 3, \ldots$ → 3Dで $\rho \propto -\zeta(-3) < 0$（引力）
\item $\lambda/4$（片端開放）: 奇数倍音 $1, 3, 5, \ldots$ → 3Dで $\rho \propto -\zeta_{\neg 2}(-3) > 0$（斥力、7倍）
\end{itemize}

PMC は要らない。片端を開けるだけ。
\end{resultbox}

%====================================================================
\part{なぜ DD と DN は根本的に違うか}
%====================================================================

\section{巻き付き数——輪ゴムの話}

\begin{storybox}[指に巻いた輪ゴム]
指に輪ゴムを巻くことを考える。

\textbf{0回巻き}: 輪ゴムは指の横にあるだけ。引っ張れば外せる。

\textbf{1回巻き}: 輪ゴムが指を1周している。引っ張っても外せない。切らない限り。

\textbf{2回巻き}: 2周。もっと外せない。

「何回巻いているか」= 巻き付き数。これは0, 1, 2, 3, ... の整数値を取る。

核心: 0回と1回の間には「連続的にいじっても越えられない壁」がある。輪ゴムを伸ばしたり縮めたりしても、巻き数は変わらない。切らないと変えられない。これが「位相的に保護されている」ということ。
\end{storybox}

\section{DD と DN は「巻き方」が違う}

\begin{mathbox}[DD = 巻いていない、DN = 巻いている]
DD（両端閉じ、$\lambda/2$）のモード: 1, 2, 3, 4, 5, ...（全整数）

全整数の世界の「形」を数学で記述すると:
$$K_1(\mathbb{Z}) = \mathbb{Z}/2$$

$\mathbb{Z}/2$ = 「0 か 1 の2択」。巻き付き数がない。
輪ゴムが指の横にあるだけの状態。

DN（片端開放、$\lambda/4$）のモード: 1, 3, 5, 7, 9, ...（奇数のみ）

奇数だけの世界の「形」:
$$K_1(\mathbb{Z}[1/2]) = \mathbb{Z}/2 \times \mathbb{Z}$$

$\mathbb{Z}$ = 「..., $-2, -1, 0, 1, 2$, ... の整数」。巻き付き数がある！
輪ゴムが指を巻いている状態。

\textbf{DD → DN に変えると、巻き付き数が「なし」から「あり」に変わる。}
\end{mathbox}

\begin{storybox}[これが意味すること]
DD の真空エネルギーは負（引力）。巻き付き数がないので、パラメータを連続的にいじれば符号を変えられる。保護がない。

DN の真空エネルギーは正（斥力）。巻き付き数があるので、パラメータを連続的にいじっても符号は変わらない。\textbf{保護されている}。

もしこの保護が物理的に本物なら:
板を離していっても DN の斥力の「正の符号」は保護される。
つまりエネルギー密度は $1/a^4$ で消えていくのではなく、
あるところでプラトーに達する（= 体積効果）。

一方、DD は保護がないので $1/a^4$ のまま消える（= 表面効果のまま）。
\end{storybox}

\section{QED はこの領域をテストしたことがない}

\begin{honestbox}[未踏領域]
物理学者は「QED が正しいから DN でも $\alpha = 5$ だ」と言う。

しかし:
\begin{itemize}
\item QED が実験で確認されたのは DD（巻き付き数なし）だけ
\item DN（巻き付き数あり）で QED をテストした人は\textbf{ゼロ}
\item 「テスト済みの領域で正しい」$\neq$「未テストの領域でも正しい」
\end{itemize}

前例:
\begin{itemize}
\item ニュートン力学は低速で完璧 → 高速で破れた（特殊相対論）
\item 古典物理はマクロで完璧 → ミクロで破れた（量子力学）
\item 両方とも「テスト済み領域では正しいが、未テスト領域で新物理が出た」
\end{itemize}

QED は DD で正しい。DN でも正しいかどうかは\textbf{実験で確かめるしかない}。

\textbf{この実験は「Boyer の予測の検証」であると同時に、
「QED を位相的に異なる真空セクターで初めてテストする」実験。}
\end{honestbox}

%====================================================================
\part{QED が「ゴミ箱に捨てたもの」}
%====================================================================

\section{繰り込み——無限大を捨てる技法}

\begin{storybox}[QED の秘密]
量子電磁力学（QED）は物理学史上最も精密な理論。電子の磁気モーメントを小数点以下12桁まで正しく予測する。

しかし QED の計算過程では「無限大」が頻繁に出現する。真空エネルギーの計算:
$$1^3 + 2^3 + 3^3 + 4^3 + \cdots = \infty$$

QED のやり方: この無限大を「繰り込み（renormalization）」で差し引く。
最終的な予測（カシミール力など）は、2つの無限大の\textbf{差}として求まり、有限で実験と合う。

\textbf{問題}: 差し引かれた無限大の中身は「ゴミ箱行き」。物理的意味がないとされてきた。
\end{storybox}

\section{ゴミ箱の中身を覗く}

\begin{resultbox}[ゼータ関数による正則化はゴミ箱の中身を見る方法]
$\zeta$ 正則化は「捨てるべき無限大」に有限の値を与える:
$$1^3 + 2^3 + 3^3 + 4^3 + \cdots = \zeta(-3) = \frac{1}{120}$$

QED はこの $1/120$ を「途中計算の数字」として使い、最終的に差を取って捨てる。

しかし $1/120 = \zeta(-3)$ には\textbf{内部構造}がある。オイラー積:
$$\zeta(-3) = \prod_{p\text{ prime}} \frac{1}{1-p^3}$$

素数 $p = 2, 3, 5, 7, \ldots$ ごとの因子に分解できる。繰り込みはこの値を丸ごと捨てるが、\textbf{素数ごとの分解が物理的に意味を持つかどうか}は一度も問われたことがない。
\end{resultbox}

\section{DN 実験 = ゴミ箱を開ける実験}

\begin{resultbox}[何が見つかるか]
DN 配置（$\lambda/4$）は $p=2$ の因子を除去する: $\zeta(-3) \to \zeta_{\neg 2}(-3)$。

これは「ゴミ箱に入っていた $\zeta(-3)$ の中身をいじる」操作。

\textbf{標準QEDの立場}: ゴミ箱の中身は物理に無関係。いじっても予測 $(\alpha = 5)$ は変わらない。

\textbf{我々の予想}: ゴミ箱の中身には素数構造があり、それは物理的。いじると $\alpha$ が変わる。

\textbf{もし $\alpha \neq 5$ が見つかったら}: 繰り込みで捨てていた真空エネルギーの算術的構造が\textbf{ずっと物理的だった}ことになる。QED の予測は正しかったが、何を捨てているかの\textbf{解釈}が不完全だった。
\end{resultbox}

\begin{storybox}[たとえ話: 家計簿]
あなたの月収が 100 万円で支出が 99 万円。手残り 1 万円。

会計士（= QED）は「手残り 1 万円」だけを報告する。月収と支出の内訳は「関係ない」と言う。

しかしもし月収の内訳（給料、ボーナス、副業、投資...）に構造があり、特定の収入源を消すと手残りの\textbf{符号が変わる}（赤字になる）としたら？

QED は「手残りの差」しか見ていない。我々は「収入の内訳（= オイラー積の素数因子）」に意味があると主張し、それを実験で確かめようとしている。
\end{storybox}

%====================================================================
\part{実験: 3つの道}
%====================================================================

\section{なぜ3つの道が必要か}

$\lambda/4$ 共振器は cm サイズで簡単に作れるが、カシミール力は nm サイズでしか測れない。このスケールの不整合を3つの方法で回避する。

\section{Route A: ナノスケール DN 構造}

\begin{expbox}[AFM で力を測る]
\textbf{アイデア}: ナノスケールで「片端開放」を実現する構造を作る。

例:
\begin{itemize}
\item 金属薄膜のスロットアレイ（穴あき板 → ノイマン条件に近い）
\item ナノコアキシャル構造（数十nmの同軸ケーブル）のオープン端
\item フォトニック結晶表面（高インピーダンス → PMC 的振る舞い）
\end{itemize}

PEC 平板とこのナノ DN 表面の間のカシミール力を AFM で測定。

\textbf{温度}: 室温（極低温不要）\\
\textbf{測定}: AFM のカンチレバーの撓み（pN レベル）\\
\textbf{判定}: 力が斥力（正）なら Boyer の予測が初確認
\end{expbox}

\section{Route B: circuit QED}

\begin{expbox}[$\lambda/4$ vs $\lambda/2$ の Lamb シフト比較]
\textbf{アイデア}: カシミール力ではなく、真空揺らぎの Lamb シフトを測る。

$\lambda/4$（奇数モード）と $\lambda/2$（全モード）の超伝導共振器に
トランスモン量子ビットを結合し、Lamb シフトの差を測定。

$\Delta(\delta f) = \sum_{n \text{ even}} \frac{g_n^2}{f_q - f_n}$

偶数モードの真空揺らぎ寄与を直接測定。

\textbf{温度}: $\sim$10 mK（希釈冷凍機）\\
\textbf{装置}: 超伝導量子ビット研究室のインフラ\\
\textbf{判定}: $\lambda/4$ と $\lambda/2$ の Lamb シフトが理論通りに異なるか
\end{expbox}

\section{Route C: MEMS}

\begin{expbox}[Si-MEMS で力 vs 距離を測定]
\textbf{アイデア}: Tang et al. (2017, Nature Photonics) の Si-MEMS プラットフォームを DN 型幾何学に適応する。

Tang の装置:
\begin{itemize}
\item コムアクチュエータで距離を制御
\item 振動ビームで力勾配を検出
\item Si ナノ構造を半導体プロセスで一括製作
\end{itemize}

T字構造の代わりにオープン端構造を配置。カシミール力を板間隔 $a$ の関数として測定し、$F \propto a^{-\alpha}$ のスケーリング指数 $\alpha$ を決定。

\textbf{温度}: 室温\\
\textbf{規模}: 100 nm スケール\\
\textbf{判定}: DN 配置で $\alpha = 5$ かつ力が正（斥力）なら Boyer 確認
\end{expbox}

%====================================================================
\part{予測}
%====================================================================

\section{3つの予測}

\begin{resultbox}[具体的な数値予測]
\textbf{予測 1}: DN 配置のカシミール力は\textbf{斥力}（正）。

\textbf{予測 2}: 3D での DN/DD エネルギー密度の比は:
$$\frac{\zeta_{\neg 2}(-3)}{\zeta(-3)} = \frac{-7/120}{1/120} = -7$$
DN 斥力は DD 引力の\textbf{7倍}の大きさ。

\textbf{予測 3}: スケーリング $F \propto a^{-5}$ は標準QEDの予測。
$\alpha \neq 5$ の逸脱が見つかれば、それは真空の位相的転移の証拠。
\end{resultbox}

%====================================================================
\part{もしうまくいったら}
%====================================================================

\section{Boyer の予測が確認されたら}

\begin{storybox}[50年越しの検証]
Boyer が1974年に予測した PEC-PMC カシミール斥力が初めて実験的に確認される。PMC なしで実現したことにより、幾何学的真空制御の新しいパラダイムが開かれる。
\end{storybox}

\section{オイラー積の物理的実在}

\begin{storybox}[素数が真空を支配する？]
$\zeta_{\neg 2}$ が物理的に実在することが確認されれば、次の疑問が自然に浮かぶ:

「$p=3$ を消したら（$\zeta_{\neg 3}$）どうなるか？」
「$p=2$ と $p=3$ を同時に消したら？」
「オイラー積の各因子が物理の何に対応するか？」

我々は既にこれらの問いに対する数論的枠組み（素数-力 対応、PFC）を持っている。
Boyer の実験的検証は、この枠組みの出発点。
\end{storybox}

\section{スケーリング逸脱が見つかったら}

\begin{storybox}[ワープへの道]
もし $\alpha \neq 5$ が見つかれば:
\begin{itemize}
\item 真空エネルギー密度が $1/a^4$ で減衰せず、定数にプラトーする
\item 「表面効果」が「体積効果」に格上げされる
\item Ford-Roman の量子エネルギー不等式をバイパスする
\item 巨視的な負のエネルギーが実現可能に
\item Alcubierre ワープ計量の物理的基盤
\end{itemize}
\end{storybox}

%====================================================================
\part{正直な評価}
%====================================================================

\begin{honestbox}[確率と価値]
\textbf{Boyer の斥力が見つかる確率}: 高い（理論的に確立、方法が正しければ）。
ただし Route A のナノ DN 構造が十分に理想的な DN 条件を実現できるかは未知。

\textbf{スケーリング逸脱が見つかる確率}: 低い（標準QEDは70年間全テストに合格）。

\textbf{実験の価値}:
\begin{itemize}
\item 斥力が見つかる（高確率）→ Boyer の50年越しの検証。確実に論文。
\item スケーリング逸脱が見つかる（低確率）→ 物理学の歴史が変わる。
\item どちらも見つからない → DN 条件の精度が不十分、または方法の限界。次へ進む。
\end{itemize}

\textbf{どちらに転んでも科学は前に進む。}
\end{honestbox}

\begin{honestbox}[この論文の独自性]
\begin{enumerate}
\item 「Boyer の計算 = オイラー積の素数2除去」という再解釈は新しい
\item 「PMC は不要、$\lambda/4$ で十分」というフレーミングは新しい
\item 3つの具体的な実験経路の提案は新しい
\item Boyer (1974) 以来50年の空白を指摘し、解決策を提示するのは新しい
\end{enumerate}

新しくないもの:
\begin{itemize}
\item Boyer の予測自体（1974年に確立）
\item ζ 正則化の物理的妥当性（カシミール実験で確認済み）
\item $\lambda/4$ 共振器の物理（高校物理）
\end{itemize}

\textbf{新しいものと古いものの組み合わせから、新しい実験が生まれる。}
\end{honestbox}

\vfill

\begin{center}
\rule{10cm}{0.4pt}\\[0.5cm]
{\Large\bfseries 1974年、Boyer は答えを出した。}\\[0.2cm]
{\Large\bfseries 2026年、我々は問いを変えた。}\\[0.2cm]
{\Large\bfseries 「PMC を見つけよ」→「片端を開けよ」。}\\[0.3cm]
{\small Wright Brothers Research Program, 2026}
\end{center}

\end{document}
